\Needspace{7\baselineskip} % REV03: subsection follows the complete electronic argument.
\subsection{Transporte lorentziano y representación electromagnética}

La sección electrónica proporciona una estructura y una coordenada
antes de su aplicación a una interacción. La realización lorentziana
de la incidencia cúbica determina, con la convención temporal declarada,
el espacio de formas en el que puede expresarse esa interacción.
Para formularla se fija una realización local orientada \((\mathcal M,g)\)
de dimensión cuatro, con fibras tangentes identificadas con el modelo
\(Q_\square\). Se especifican además una conexión \(\nabla\) con
\(\nabla g=0\), una uno-forma potencial \(A\in\Omega^1(\mathcal M)\),
su curvatura \(F=dA\) y una representación de carga de la sección.
Estos datos son los argumentos de la representación física que sigue.

La dualidad de Hodge satisface \(\star^2=-I\) sobre las dos-formas.
La separación eléctrica y magnética corresponde a una elección temporal
en la descomposición \(1+3\); \(F\) y \(\star F\) conservan su
relación dentro del mismo espacio de dos-formas. La neutralidad de una
representación de carga significa \(q=0\), mientras la sección puede
conservar orientación, fase y transporte interno. La ausencia de
acoplamiento eléctrico en esa representación tiene un contenido distinto
de la ausencia de información de fase.

En esta realización dimensional, sean \(m_e>0\), \(q\) la carga,
\(\gamma(s)\) una trayectoria parametrizada y \(u=\dot\gamma\)
su cuadrivelocidad. La fuerza de Lorentz se expresa por
\begin{equation}
 m_e\nabla_u u^\mu=qF^\mu{}_{\nu}u^\nu.
 \label{eq:lorentz-electron}
\end{equation}
La antisimetría de \(F\) implica
\(F_{\mu\nu}u^\mu u^\nu=0\). Por consiguiente,
\[
 \frac{d}{ds}g(u,u)=2g(\nabla_u u,u)=0.
\]
Se conserva la normalización de la trayectoria siempre que \(\nabla\)
sea métrica y se cumpla \eqref{eq:lorentz-electron}. La masa utilizada
es la realización dimensional de la sección electrónica precedente;
el parámetro de acoplamiento se identifica en la convención electromagnética
adoptada. La deducción concreta del potencial y de una configuración
de campo constituye un problema adicional con condiciones de frontera
propias. El resultado aquí expuesto determina el dominio, los datos
de la representación y la conservación que su ecuación implica.

La inversión de rutas del TPK y la orientación de una línea electrónica
ofrecen entonces una comparación precisa: cada una tiene su involución,
su transporte y su representación. El antecedente de Wheeler y Feynman
sitúa históricamente la pregunta; la equivalencia de dos realizaciones
se decide mediante sus aplicaciones y condiciones de compatibilidad.
Esta disposición permite mantener unido el origen aritmético del electrón
con la formulación de su interacción, haciendo explícita la función
de cada etapa.
